\clearpage
\VerdierPhysicalPageAnchor{0234}%
\begingroup\small\setlength{\parskip}{0pt}\setlength{\emergencystretch}{1.5em}
\setlength{\abovedisplayskip}{4pt plus 2pt minus 2pt}%
\setlength{\belowdisplayskip}{4pt plus 2pt minus 2pt}%

\VerdierSemanticAnchor{ch4d-4-1-9-cont}
\noindent $\bigl({}^{1}\!Y(p,q),{}^{1}\!\delta(p,q,r)\bigr)$（それぞれ
$\bigl({}^{2}\!Y(p,q),{}^{2}\!\delta(p,q,r)\bigr)$）を、複体 $X$ を第一次数
（それぞれ第二次数）でフィルター付けして定まる、$K(\mathcal A)$ に値をとる
スペクトル対象 (4.1.7) とする。定義から直ちに次の関係を得る：

\VerdierSemanticAnchor{ch4d-4-1-9-1}
\begin{equation}
\left\{
\begin{aligned}
{}^{1}\!Y(q-1,q)&=X^{-q,\bullet}[q],\\
{}^{2}\!Y(q-1,q)&=X^{\bullet,-q}[q],\\
{}^{1}\!\delta(q-1,q,q+1)&=-\bigl(d_1^{-q-1,\bullet}\bigr)[q+1],\\
{}^{2}\!\delta(q-1,q,q+1)&=-\bigl(d_2^{\bullet,-q-1}\bigr)[q+1].
\end{aligned}
\right.
\tag{4.1.9.1}
\end{equation}
これらの公式は (4.6) で用いられる。

\VerdierSemanticAnchor{ch4d-4-2}
\medskip\noindent\textbf{4.2. $\operatorname{comp}(\mathcal A)$ の
フィルター付き対象に付随する、$D(\mathcal A)$ に値をとるスペクトル対象。}

\VerdierSemanticAnchor{ch4d-4-2-scope}
\begin{center}\itshape
この項および本節の以下の各項では、$\mathcal A$ はアーベル圏を表す。
\end{center}

\VerdierSemanticAnchor{ch4d-4-2-1}
\noindent\textbf{4.2.1.}\quad
$X$ を $\operatorname{comp}(\mathcal A)$ の対象とし、
$X(q)$、$q\in\mathbb Z$ を複体 $X$ の
$\widetilde{\mathbb Z}$ 型増大フィルトレーションとする：

\VerdierSemanticAnchor{ch4d-4-2-1-1}
\begin{equation}
0=X(-\infty)\subset\cdots\subset X(q)\subset X(q+1)
\subset\cdots\subset X(+\infty)=X.
\tag{4.2.1.1}
\end{equation}

\VerdierSemanticAnchor{ch4d-4-2-1-associated-object}
(4.1.1) で想起した構成により、$K(\mathcal A)$ に値をとるスペクトル対象
$\bigl(\mathcal X(p,q),\delta(p,q,r)\bigr)$ を得る。さらに標準関手
$K(\mathcal A)\longrightarrow D(\mathcal A)$ を適用すると、
$D(\mathcal A)$ に値をとるスペクトル対象を得るが、これも記号の濫用により
$\bigl(\mathcal X(p,q),\delta(p,q,r)\bigr)$ と書く。以下では、このスペクトル対象と
同型な、$D(\mathcal A)$ に値をとるスペクトル対象を記述する。

\VerdierSemanticAnchor{ch4d-4-2-2}
\noindent\textbf{4.2.2.}\quad
そこで、
$0=X(-\infty)\subset\cdots\subset X(q)\subset\cdots\subset
X(+\infty)=X$ を $\operatorname{comp}(\mathcal A)$ の
フィルター付き対象とする。次のようにおく：

\VerdierSemanticAnchor{ch4d-4-2-2-1}
\begin{equation}
X(p,q)=X(q)/X(p),
\tag{4.2.2.1}
\end{equation}
ここで $(p,q)$ は $\mathcal{F}\ell_1(\widetilde{\mathbb Z})$ の任意の元である。
（ここで商は、もちろんアーベル圏 $\operatorname{comp}(\mathcal A)$ における
意味でとる。）次を

\VerdierSemanticAnchor{ch4d-4-2-2-2}
\begin{equation}
\varphi\bigl((p,q),(p',q')\bigr)\colon
X(p,q)\longrightarrow X(p',q'),
\tag{4.2.2.2}
\end{equation}
で表す。
\endgroup
